\section{Análisis de riesgos y plan de respuesta a riesgos}
\subsection{Amenazas y Vulnerabilidades (El entorno hostil)}
\subsubsection{1. Amenazas Ambientales (Entorno Físico)}
\begin{itemize}
    \item ¿Las computadoras están expuestas a polvo, rebaba de metal, calor extremo o lluvia? Polvo y calor extremo.
    \item Energía: Al encender maquinaria pesada (compactadoras, grúas), ¿notan bajones de luz en las computadoras? (Esto daña discos duros). Casi todos los días, por eso se cuenta con un generador y electricista de respaldo.
\end{itemize}
\subsubsection{Amenazas Humanas (Fraude Interno)}
\begin{itemize}
    \item ¿Es posible que un operador modifique el peso en el sistema manualmente para pagar de más a un proveedor "amigo"? No.
    \item ¿Comparten usuarios y contraseñas entre varios empleados de la báscula? Sí.
\end{itemize}
\subsubsection{2. Vulnerabilidades Técnicas}
\begin{itemize}
    \item ¿El software de la báscula es antiguo (ej. corre en Windows 7 o XP)? Sí.
    \item ¿Si se va la luz o el internet, pueden seguir operando manualmente o se detiene la compra? Se detiene totalmente.
\end{itemize}

\subsection{Evaluación de Riesgo (Impacto y Probabilidad)}
\subsubsection{1. Impacto en la Tríada CID}
\begin{itemize}
    \item \textbf{Integridad (Crítico):} Si alguien altera la base de datos de pesajes, ¿se dan cuenta inmediatamente o hasta el cierre de mes? Después del corte de semana.
    \item \textbf{Disponibilidad:} Si el sistema de facturación falla un viernes de pago, ¿qué pasa con los proveedores? Se habla y se retrasa el pago.
    \item \textbf{Confidencialidad:} Si la competencia supiera sus precios de compra o su lista de grandes proveedores industriales, ¿perderían negocio? No deben de conocerlos en la competencia.
\end{itemize}
\subsubsection{2. Probabilidad}
\begin{itemize}
    \item ¿Con qué frecuencia fallan los equipos por polvo o suciedad? Rara vez, pero ha pasado.
    \item ¿Qué tan seguido reciben correos falsos de "facturas" (phishing) en el área administrativa? No casi no, o normalmente sólo ignoran los correos que no conocen.
\end{itemize}

\subsection{Controles y Tratamiento (Soluciones)}
\subsubsection{1. Controles Actuales}
\begin{itemize}
    \item ¿Tienen UPS (No-Breaks) de grado industrial para las computadoras de operación? No.
    \item ¿Se hacen respaldos de la base de datos de la báscula? ¿En dónde (USB, Nube, Disco externo)? En usb y nube, se considera empezar a usar disco duro externo este mes.
\end{itemize}
\subsubsection{2. Presupuesto y Viabilidad}
\begin{itemize}
    \item ¿Estarían dispuestos a invertir en equipos de "uso rudo" o prefieren reemplazar equipos baratos frecuentemente? Consideramos comprar máquinas de uso rudo.
    \item ¿Tienen presupuesto para licencias de software original o trabajan con lo que venía en la PC? Con lo que venía la pc.
\end{itemize}

\subsection{Tips para tu Reporte (Caso Reciclaje)}
\begin{itemize}
    \item \textbf{Riesgo Top 1 (Integridad):} Modificación no autorizada de pesos/precios. Tratamiento: Controles de acceso lógico (usuarios únicos) y logs de auditoría. ¿Conocen este riesgo? Sí.
    \item \textbf{Riesgo Top 2 (Disponibilidad):} Falla de hardware por entorno sucio/eléctrico. Tratamiento: Mantenimiento preventivo, UPS, ubicación física protegida. ¿Están protegidos? No están tan protegidos de este riesgo.
    \item \textbf{Riesgo Top 3 (Legal/Confidencialidad):} Robo de datos de proveedores industriales. Tratamiento: Cifrado de discos o contratos de confidencialidad (NDAs). ¿Tienen cifrado? No, contrato de confidencialidad sí hay.
\end{itemize}

\subsection{Matrices de Riesgo}

\begin{itemize}
    \item \textbf{PC Operativa (báscula)}
    \begin{itemize}
        \item \textbf{Amenaza:} Fuente de luz inestable
        \item \textbf{Vulnerabilidad:} Inexistencia de un No-brake
        \item \textbf{Riesgo:} Daño físico al disco duro
        \item \textbf{Tratamiento:} Adquisición de un No-brake
    \end{itemize}

    \item \textbf{Sistema Operativo}
    \begin{itemize}
        \item \textbf{Amenaza:} Falla en el software
        \item \textbf{Vulnerabilidad:} Sistema operativo obsoleto
        \item \textbf{Riesgo:} Fallos y brechas de seguridad y paro total de operaciones
        \item \textbf{Tratamiento:} Adquisición de equipo más moderno e inversión en licencias para software empresarial
    \end{itemize}

    \item \textbf{Datos de pesaje}
    \begin{itemize}
        \item \textbf{Amenaza:} Fraude interno / Modificación no autorizada
        \item \textbf{Vulnerabilidad:} Cuentas y contraseñas compartidas
        \item \textbf{Riesgo:} Integridad de la información
        \item \textbf{Tratamiento:} Implementación de usuarios únicos y activación de logs
    \end{itemize}

    \item \textbf{PC Operativa}
    \begin{itemize}
        \item \textbf{Amenaza:} Condiciones hostiles
        \item \textbf{Vulnerabilidad:} Uso de software de base
        \item \textbf{Riesgo:} Pérdida de confidencialidad
        \item \textbf{Tratamiento:} Instalación de antivirus y restringir puertos USB
    \end{itemize}

    \item \textbf{Red / Servidor}
    \begin{itemize}
        \item \textbf{Amenaza:} Ransomware
        \item \textbf{Vulnerabilidad:} No existe uso de software especializado de seguridad
        \item \textbf{Riesgo:} Cifrado de archivos esenciales
        \item \textbf{Tratamiento:} Instalación de antivirus
    \end{itemize}

    \item \textbf{Información fiscal y de proveedor}
    \begin{itemize}
        \item \textbf{Amenaza:} Robo de información
        \item \textbf{Vulnerabilidad:} Archivos Excel sin cifrado en disco local o USB
        \item \textbf{Riesgo:} Fuga de datos sensibles
        \item \textbf{Tratamiento:} Dejar de usar canales públicos para acceder a estos datos
    \end{itemize}
\end{itemize}


\subsection{Costos de mejora}
\begin{itemize}
    \item No-Break CyberPower 1500VA: \$5,558.8
    \item Windows 11 Pro (OEM): \$2,949.00
    \item Antivirus (Anual): Eset NOD32 (5 Usuarios): \$889.00
    \item Kit de Mantenimiento: \$300.0
    \item Mano de Obra: \$800.0
    \item \textbf{Estimado total: \$10,496.}
\end{itemize}
